\subsection{CLUSTER POINT OF A FILTER BASE}

DEFINITION 2. In a topological space X, a point x is a cluster point of a filter base B on X if it lies in the closure of all the sets of B.

If \(x\) is a cluster point of a filter base \(\mathfrak{B}\) , it is also a cluster point of every equivalent filter base by reason of § 6, no. 3, corollary to Proposition 4; in particular, \(x\) is a cluster point of the filter whose base is \(\mathfrak{B}\) .

PROPOSITION 3. A point x is a cluster point of a filter base B if and only if every set of a fundamental system of neighbourhoods of x meets every set of B.

This follows immediately from the definitions.

This proposition and Corollary 2 to Proposition 1 of § 6, no. 2 show that the property ``x is a cluster point of the filter F'' is equivalent to the property ``there is a filter which is finer than both F and the neighbourhood filter of x''. In other words:

PROPOSITION 4. A point \(x\) is a cluster point of a filter \(\mathfrak{F}\) if and only if there is a filter finer than \(\mathfrak{F}\) which converges to \(x\) .

In particular, every limit point of a filter \(\mathfrak{F}\) is a cluster point of \(\mathfrak{F}\) .

COROLLARY. An ultrafilter U converges to a point x if and only if x is a cluster point of U.

If x is a cluster point of a filter F, it is also a cluster point of every filter coarser than F; likewise, if we replace the topology of X by a coarser topology, x remains a cluster point of F in the new topology.

The set of cluster points of a filter base \(\mathfrak{B}\) on \(\mathbf{X}\) is by definition the set \(\bigcap_{\mathbf{M} \in \mathfrak{B}} \overline{\mathbf{M}}\) , whence

PROPOSITION 5. The set of cluster points of a filter base on a topological space X is closed in X.

PROPOSITION 6. Let \(\mathfrak{B}\) be a filter base on a subset A of a topological space X. Then every cluster point of \(\mathfrak{B}\) in X belongs to \(\overline{\mathbf{A}}\) ; and conversely every point of \(\overline{\mathbf{A}}\) is a limit point of a filter on A.

The first assertion is trivial; on the other hand, if \(x \in \overline{A}\) , the trace on A of the neighbourhood filter of x in X is a filter on A which evidently converges to x.

Remark. A filter on a topological space need have no cluster points (and a fortiori no limit points); for example, in an infinite discrete space the filter of complements of finite subsets has no cluster points. Spaces in which every filter has a cluster point play an important role in mathematics, and we shall study them in § 9.

